%%%%%%%%%%%%%%%%%%%%%%%%%%%%%%%%%%%%%%%%%%%%%%%%%%%%%%%%%%%%%%%%%%%%%%%%%%%%%%%%%%%%%%%%%%%%%%%%%%%%%%
%
%   Filename    : chapter_2.tex 
%
%   Description : This file will contain your review of related works.
%                 
%%%%%%%%%%%%%%%%%%%%%%%%%%%%%%%%%%%%%%%%%%%%%%%%%%%%%%%%%%%%%%%%%%%%%%%%%%%%%%%%%%%%%%%%%%%%%%%%%%%%%

\chapter{Related Works}
\label{sec:relatedworks}

\section{Dense Small-Object Detection}

Dense small-object detection \cite{smallObject1,smallObject2} has been studied across entomology, agriculture, and digital pathology, where automation replaces labor-intensive manual quantification. Following the COCO convention \cite{Lin2014COCO}, an object is considered small when it occupies fewer than $32 \times 32$ pixels. In mosquito surveillance, EggCountAI \cite{Javed2023} uses Mask R-CNN to achieve high egg-counting accuracy under controlled imaging. Although accurate, the model's large parameter count limits real-time, low-power deployment. Subsequent work extended detection to field-collected ovitrap images with CNN-based frameworks \cite{Hamesse2023,NaranjoAlcazar2024,vicente2024aedesaegyptieggcounting}. These approaches enable rapid quantification of \textit{Aedes} eggs, a crucial step for predicting and mitigating mosquito-borne disease outbreaks.

In agricultural phenotyping, YOLO-based pipelines are common for seedling and grain counting from drone \cite{Chen2024} and smartphone imagery \cite{Santiago2024}. Automating dense agricultural target estimation accelerates yield forecasting and breeding, which were previously bottlenecked by manual counting. Large-scale benchmarks such as the Sorghum-Grain-Count (SGC) dataset \cite{Bakshi2025} provide thousands of densely packed grains per panicle, demanding high-capacity detectors that can handle severe occlusions. Reconstruction-based methods \cite{Freeman2023} further improve accuracy by exploiting 3D spatial information.

\begin{sloppypar}
Computational pathology benefits from automated detection of histologic structures to support celiac disease diagnosis \cite{Stoleru2022,Sali2019}. Frameworks like DeGPR \cite{Tyagi2023} improve multi-class cell detection and counting in crowded microscopic images through deep guided posterior regularisation. Despite diverse imaging modalities, all domains share the challenge of densely packed small objects that strain detector localisation.
\end{sloppypar}

\begin{sloppypar}
Modern detectors address small-object sensitivity through multi-scale feature aggregation and architectural optimisation. YOLO families (YOLOv4 \cite{Bochkovskiy2020YOLOv4}, YOLOv8 \cite{yolov8}) separate backbone, neck, and head, using feature pyramids to improve localisation while retaining real-time speed. Transformer-based detectors such as DETR \cite{Carion2020DETR} treat detection as set prediction, removing anchors and non-maximum suppression. Real-time variants like RT-DETR \cite{Zhao2023RTDETR} match or exceed YOLO performance on real-time COCO benchmarks, offering alternative trade-offs between backbone extraction, global attention, and prediction heads. These differences motivate module-specific analysis for efficiency.
\end{sloppypar}

\section{Model Compression Strategies}

\begin{sloppypar}
Although architectural advances have boosted accuracy and speed, most domain-specific counting systems prioritise performance over parameter efficiency. Heavyweight YOLO variants, Mask R-CNN backbones, and transformer-based detectors typically run on GPU workstations, limiting field or edge use \cite{He2025}. Efficiency-oriented small-object detectors such as QueryDet \cite{Yang2022QueryDet} reduce the cost of high-resolution feature processing through cascaded sparse queries, which is complementary to compressing the network itself. Model-compression strategies \cite{He2025}---including pruning, quantisation, and lightweight redesign---aim to reduce parameter count and computational cost. Structured filter pruning \cite{Li2017FilterPruning,Luo2017ThiNet} removes entire convolutional channels based on importance criteria, yielding hardware-friendly reductions that are typically paired with fine-tuning to recover accuracy. In detection, anchor-level pruning \cite{Bonnaerens2022AnchorPruning} improves efficiency for anchor-based detectors. However, pruning may degrade multi-scale feature aggregation modules, such as YOLO necks and transformer encoders, which are essential for small-object localisation; this thesis examines this module by module.
\end{sloppypar}

To counteract pruning-induced degradation, Knowledge Distillation (KD) can transfer the original model's knowledge to the pruned one. Beyer et al. \cite{beyer2022knowledge} show that distillation is most effective when framed as function matching: teacher and student receive identical augmented inputs, and training runs over very long schedules. For detection, Focal and Global Distillation (FGD) \cite{yang2022focal} addresses the foreground--background imbalance in detector features by separating the two and adding global relation distillation, while Channel-Wise Distillation (CWD) \cite{shu2021channel} aligns the channel-wise activation distributions of teacher and student for dense prediction. ScaleKD \cite{zhu2023scalekd} targets small-object detection by decoupling features by scale and using a cross-scale assistant to refine noisy teacher predictions before distillation. These methods were developed for general or small-object detectors rather than for pruned students; this thesis examines their use in recovering pruned detectors. In specialised tasks, KD alone has been applied to fast nuclei instance segmentation in FKDNuSeg \cite{Zhao2026_FKDNuSeg}.

\begin{sloppypar}
Efficient edge inference also benefits from reduced-precision quantisation. Standard Post-Training Quantisation (PTQ) may erase weak activation signals, which can harm small-object recall. Specialised quantisation methods target detector-specific failure modes. Reg-PTQ \cite{ding2024regptq} targets the regression branch of detectors, whose Laplace-like weight distributions are poorly matched to uniform quantisation bins; it uses a learnable logarithmic-affine quantiser and calibrates on the global detection loss, without full QAT. For transformer-based models, SmoothQuant \cite{xiao2023smoothquant}, originally proposed for large language models, shifts quantisation difficulty from activations to weights, enabling INT8 quantisation. Activations in transformers contain large outliers concentrated in a few channels, which makes them difficult to quantise to INT8, whereas weights are comparatively easy to quantise. SmoothQuant divides each activation channel by a scaling factor and multiplies the corresponding weights by the same factor. Consequently, the layer output remains mathematically unchanged, but activation ranges become smoother and weights absorb the quantisation difficulty. Q-DETR \cite{xu2023qdetr} addresses distribution distortion in the queries of low-bit DETR through distribution rectification distillation.
\end{sloppypar}

\begin{sloppypar}
Several studies combine two of these compression techniques. For pruning and distillation, PKD-YOLOv8 \cite{Yu2025_PKD} couples filter pruning with knowledge distillation for lightweight pest detection. For quantisation and distillation, Wei et al. \cite{Wei2018} propose Quantization Mimic to train tiny detection networks such as channel-reduced VGG and ResNet. Their method first quantises the large teacher network to restrict the search space, and then trains a quantised small student network to match the teacher's feature maps. For pruning and quantisation, Endrawati et al. \cite{Endrawati2021} apply iterative layer pruning and integer quantisation to YOLOv3-Tiny for edge deployment. They prune selected convolutional layers with repeated retraining on Pascal VOC person detection, reducing weight storage while preserving detection accuracy.
\end{sloppypar}

Despite these advances, the works reviewed above combine at most two of the three techniques: pruning with KD \cite{Yu2025_PKD}, quantisation with distillation \cite{Wei2018}, or pruning with quantisation \cite{Endrawati2021}. Compression research is also typically evaluated on generic benchmarks rather than dense small-object counting. A unified, multi-step pipeline---architecture-aware pruning followed by distillation and specialised quantisation---remains underexplored for dense small-object counting, and systematic studies of how such pipelines affect distinct modules across heterogeneous, high-density datasets are lacking. Addressing this gap motivates a comprehensive evaluation of a unified compression framework across multiple modules, detection architectures, and extreme-density datasets to establish practical guidelines for parameter-efficient edge deployment.

\begin{comment}
% IPR acknowledgement: the contents within this comment are from Ethel Ong's slides on RRL.
% Guide on Writing your Related Works chapter
%   1. Identify the keywords with respect to your research
%   2. Find references using these keywords
%   3. For each of the references that you find, check relevance
%   4. Identify a set of criteria for comparison
%   5. Write a summary focusing on what, how, findings, why
%   6. At the end of each section, show a Table of Comparison of the related works and your proposed project/system
\end{comment}
